\documentclass[10pt]{article}
\usepackage[margin=2cm]{geometry}
\usepackage{booktabs}
\usepackage{graphicx}
\usepackage{hyperref}
\usepackage{amsmath}

\begin{document}

\title{CSCI 5942: AI Engineering Assignment 5}
\author{Bill Hsu}
\maketitle

% =====================================================================
\section*{Setup}

For this assignment, all GPU runs used one NVIDIA A40 (44\,GiB) 
on an NCSA Delta
\texttt{gpuA40x4} node with the \texttt{pytorch-conda/2.12} module.
Weights, gradients and AdamW moments are fp32; compute runs in bf16
under autocast. Every graded run is \texttt{gpt2-medium}, one epoch
(2{,}629 steps, batch 16), with the given configs. Evaluation is
\texttt{eval.py}: greedy decoding on the 630 test meaning
representations, corpus BLEU (sacrebleu) against all references, slot
coverage, and teacher-forced loss on 200 FineWeb-Edu blocks. All tests
pass: \texttt{tests/test\_mask.py} (3), \texttt{tests/test\_lora.py} (3)
and \texttt{tests/test\_pretrained.py} (1).

% =====================================================================
\section*{Part 1: The Memory Model}

Predicted bytes are 16 per trainable weight (fp32 weight, gradient, two
Adam moments) and 4 per frozen weight. Our model class does not tie
\texttt{wte} and \texttt{lm\_head}, so $N$ includes the embedding twice
(e.g.\ 406.3M for \texttt{gpt2-medium}, not 354.8M). A block has
$12d^2 + 13d$ parameters; LoRA on the four Linear layers adds
$r(4d + 2d + 5d + 5d) = 16rd$ per block. Measured rows are 50-step runs.

\begin{table}[htbp]
\caption{Predicted model state against the measured peak
(\texttt{torch.cuda.max\_memory\_allocated}). Residual = peak $-$ predicted.}
\label{tab:mem}
\centering
\begin{tabular}{lrrrr}
\toprule
Model & Trainable & Predicted (params + grads + opt.) & Measured peak & Residual \\
\midrule
gpt2, full                & 163{,}037{,}184   & 2.43\,GiB  & 3.83\,GiB & 1.40\,GiB \\
gpt2-medium, full         & 406{,}286{,}336   & 6.05\,GiB  & 8.63\,GiB & 2.57\,GiB \\
gpt2-medium, LoRA $r=4$   & 1{,}572{,}864     & 1.54\,GiB  & 4.25\,GiB & 2.71\,GiB \\
gpt2-medium, LoRA $r=64$  & 25{,}165{,}824    & 1.89\,GiB  & 4.57\,GiB & 2.68\,GiB \\
gpt2-xl, full             & 1{,}638{,}022{,}400 & 24.41\,GiB & (predict only) & \\
gpt2-xl, LoRA $r=4$       & 4{,}915{,}200     & 6.18\,GiB  & (predict only) & \\
\bottomrule
\end{tabular}
\end{table}

\begin{itemize}
\item \textbf{Where the residual comes from.} The residual is mostly
  activations saved for the backward pass (plus the
  $B \times T \times 50257$ logits and autocast's bf16 copies of the
  weights). It scales with batch, sequence length, width and depth, but
  not with the number of trainable parameters.
\item \textbf{Why LoRA barely changes it.} The residual stays almost the
  same under LoRA (2.57 vs.\ 2.71\,GiB) because adapters sit in every
  block, so the backward pass still runs through the whole network and
  keeps every layer's activations. LoRA removes only the gradients and
  Adam state of the frozen weights.
\item \textbf{\texttt{gpt2-xl}, full fine-tuning.} Based on our calculations,
  it needs 24.4\,GiB of model state before any activations. That alone
  is more than a 24\,GiB card holds. If the residual grows with
  depth $\times$ width (about $3\times$ medium's, so roughly 8\,GiB),
  the total is about 32\,GiB, which fits on a 40\,GiB A40 with little
  room to spare. This residual is an estimate, not a measurement.
\item \textbf{\texttt{gpt2-xl}, LoRA $r=4$.} It needs 6.2\,GiB plus
  activations (about 14\,GiB) and fits on either card.
\end{itemize}

% =====================================================================
\section*{Part 2: The Template and the Mask}

\texttt{sft.describe} for one training pair (left: the input token
$x_t$; right: whether predicting $x_{t+1}$ is scored):

\begin{verbatim}
          name  -
             [  -
            Al  -
         iment  -
            um  -
            ],  -
          area  -
             [  -
          city  -
        centre  -
             ]  -
            \n  -
   Description  -
             :  scored
            Al  scored
         iment  scored
            um  scored
            is  scored
            in  scored
           the  scored
          city  scored
        centre  scored
             .  scored
\end{verbatim}

The first scored position is the colon, whose next token is the first
target token \texttt{" Al"}; the last is the period, whose next token is
\texttt{<|endoftext|>}. All three tests in \texttt{tests/test\_mask.py}
pass.

Prompts average 34.8 tokens and targets 25.0 on the training set, so
scoring the prompt would put 58\% of the loss on reproducing
meaning-representation syntax, which the model never has to generate at
test time, and dilute the signal for the description.

% =====================================================================
\section*{Part 3: Full Fine-Tuning}

\begin{table}[htbp]
\caption{The untouched model and one epoch of full fine-tuning.
FineWeb is teacher-forced loss on held-out FineWeb-Edu (nats/token).}
\label{tab:fft}
\centering
\begin{tabular}{lrrrrrr}
\toprule
Model & BLEU & Slot coverage & Dev loss & FineWeb loss & Peak memory & Train time \\
\midrule
gpt2-medium, base & 9.5  & 33.9\% & 2.837$^{*}$ & 3.017 & --        & --    \\
gpt2-medium, full & 63.2 & 97.3\% & 0.994       & 3.826 & 9.26\,GiB & 335\,s \\
\bottomrule
\end{tabular}

\smallskip
\footnotesize $^{*}$\texttt{train.dev\_loss} on the same 512 dev pairs,
run on CPU in fp32 because \texttt{train.py} writes no summary for the
untrained model.
\end{table}

The same meaning representation before and after:

\begin{quote}\small
\texttt{name[Blue Spice], eatType[coffee shop], area[city centre]}\\[2pt]
\textbf{Base:} \verb|Description:\nDescription:\nDescription:|\ldots
(repeats until the length limit)\\
\textbf{Full:} Blue Spice is a coffee shop located in the city centre.
\end{quote}

The FineWeb loss rose from 3.017 to 3.826 nats per token (+0.81, or
27\%): one epoch of updating every weight on restaurant descriptions
measurably damaged the model's general language modeling.

The full run's peak (9.26\,GiB) is higher than Part 1's 50-step peak
(8.63\,GiB) because a full epoch eventually meets the longest batches,
and activations grow with sequence length.

% =====================================================================
\newpage

\section*{Part 4: LoRA and the Rank Knee}

\begin{figure}[htbp]
\centering
\includegraphics[width=\linewidth]{knee.png}
\caption{Dev loss (left) and test BLEU (right) against trainable
parameters. Dashed: full fine-tuning. The BLEU axis spans under
2 points.}
\label{fig:knee}
\end{figure}

\begin{table}[htbp]
\caption{LoRA ($\alpha = 2r$, lr $2\times10^{-4}$, all four Linear
layers) against full fine-tuning (lr $5\times10^{-5}$). Base FineWeb
loss is 3.017.}
\label{tab:knee}
\centering
\begin{tabular}{lrrrrrr}
\toprule
Run & Trainable & Dev loss & BLEU & FineWeb loss & Peak memory & Train time \\
\midrule
LoRA $r=1$  & 393{,}216       & 1.125 & 61.4 & 3.177 & 4.96\,GiB & 175\,s \\
LoRA $r=4$  & 1{,}572{,}864   & 1.087 & 61.6 & 3.136 & 4.97\,GiB & 184\,s \\
LoRA $r=16$ & 6{,}291{,}456   & 1.050 & 61.9 & 3.160 & 5.04\,GiB & 185\,s \\
LoRA $r=64$ & 25{,}165{,}824  & 1.021 & 63.0 & 3.174 & 5.30\,GiB & 194\,s \\
Full        & 406{,}286{,}336 & 0.994 & 63.2 & 3.826 & 9.26\,GiB & 335\,s \\
\bottomrule
\end{tabular}
\end{table}

\begin{itemize}
\item \textbf{Where the knee is.} There is no sharp knee in this range.
  Dev loss falls steadily, by about 0.03 to 0.04 nats for each
  $4\times$ in rank, and has not flattened by $r=64$. BLEU is nearly
  flat: $r=1$, with 0.1\% of the weights, already reaches 61.4 of full
  fine-tuning's 63.2, so for the task metric the knee is at or below
  $r=1$.
\item \textbf{Does LoRA reach full fine-tuning?} On BLEU, yes: $r=64$
  comes within 0.2 points of full fine-tuning, which is inside what one
  seed can resolve. On dev loss, no: LoRA does not reach full
  fine-tuning (1.021 vs.\ 0.994 at $r=64$).
\item \textbf{Which method forgot more.} Full fine-tuning forgot far
  more: its FineWeb loss rose by 0.81 nats, against 0.12 to 0.16 for
  every LoRA rank, with no clear trend in rank.
\item \textbf{Cost.} LoRA also used about half the peak memory and
  trained in about 55\% of the time, since the frozen weights need no
  weight-gradient matmuls and no optimizer update.
\item \textbf{Comparison with the paper.} Unlike Table~3 of the LoRA
  paper (Hu et al., 2021; 70.4 vs.\ 68.2), LoRA does not beat full
  fine-tuning here. Our setting differs: one epoch instead of five,
  greedy decoding instead of beam search, sacrebleu instead of the
  official scorer, and adapters on all four Linear layers at $r=1$ to
  $64$, where the paper adapts only $W_q$ and $W_v$ at $r=4$. At the
  same 0.39M budget as the paper's configuration, our LoRA $r=1$
  reaches 61.4 BLEU.
\end{itemize}
% =====================================================================
\newpage
\section*{Part 5: Which Matrices}

\begin{table}[htbp]
\caption{About 1.5M trainable parameters spent three ways.}
\label{tab:budget}
\centering
\begin{tabular}{lrrr}
\toprule
Config & Trainable & Dev loss & BLEU \\
\midrule
\texttt{lora-r4} (all four, $r=4$)            & 1{,}572{,}864 & 1.087 & 61.6 \\
\texttt{budget-mlp-r12} (\texttt{c\_fc}, $r=12$)  & 1{,}474{,}560 & 1.091 & 61.8 \\
\texttt{budget-attn-r16} (\texttt{c\_attn}, $r=16$) & 1{,}572{,}864 & 1.126 & 62.0 \\
\bottomrule
\end{tabular}
\end{table}

Based on the results, we can conclude:
\begin{itemize}
\item \textbf{Where the budget should go.} By dev loss, a fixed budget
  is best spread across all four matrices at low rank (1.087), with
  MLP-only close behind (1.091) and attention-only clearly worst
  (1.126).
\item \textbf{BLEU cannot rank them.} The three BLEU scores span only
  0.4 points, which one seed per configuration cannot resolve.
\item \textbf{Agreement with the paper.} This agrees with Table~5 of
  the LoRA paper (Hu et al., 2021), where adapting more weight types at
  lower rank beats one type at higher rank: there, $W_q, W_v$ or all
  four attention matrices beat $W_q$ alone.
\item \textbf{Caveat.} The paper compares only attention matrices on
  GPT-3 and does not test the MLP, so the agreement is in direction,
  not in numbers.
\end{itemize}

% =====================================================================
\section*{When Is Updating Every Weight Worth It?}

Full fine-tuning is worth it when the last bit of in-domain quality
matters, the model is small enough that 16 bytes per weight fits on the
hardware, and the model will serve only this one task. On this task,
full fine-tuning buys little. Full fine-tuning reached the best dev loss
(0.994 against LoRA $r=64$'s 1.021) but essentially the same BLEU
(63.2 vs.\ 63.0, and 61.4 already at $r=1$; Figure~\ref{fig:knee}). It
cost about twice the peak memory and time (Table~\ref{tab:mem}:
6.05\,GiB of state against 1.54\,GiB; 9.26 vs.\ about 5\,GiB peak), and
it forgot five to seven times as much (FineWeb +0.81 nats against
+0.12 to +0.16). LoRA also stores each task as a few megabytes of
adapters on one shared base model, while full fine-tuning stores a new
1.6\,GB copy per task.

For a model ten times this size (about 4B parameters), memory changes
the answer the most. Full fine-tuning needs about 65\,GB of model state,
more than one A40 holds, so it needs the sharding from Assignment~4
before any activations. LoRA needs about 16\,GB of frozen weights and
fits on one GPU. The quality gap measured by the knee would likely
shrink, not grow, at that size (the paper reports LoRA matching or
beating full fine-tuning on GPT-3), but that comes from the paper, not
from our runs. Forgetting stays a reason to prefer LoRA at any size,
because it leaves the base weights untouched.

\end{document}
